% [NERD-PROTOCOL] CCIS (llncs.cls) skeleton — hypothesis-first structure per the
% hypothesis package. ANONYMIZED for double-blind review.
% Page budget (regular, 12-16pp): S1 1.5pp / S2 2pp / S3 2.5pp / S4 3.5pp / S5 2pp /
% S6 1pp / refs 1.5pp — tracked here, adjusted at day-13 writing start.
% NO citations are entered until verified against the source (citation-integrity rule);
% \citeneed{} marks every spot that needs one.
\documentclass{llncs}
\usepackage{graphicx}
\usepackage{booktabs}
\newcommand{\citeneed}[1]{\textcolor{red}{[citation needed: #1]}}
\usepackage{xcolor}

\begin{document}

% TITLE POLICY: deferred until results exist (hypothesis package section 5).
% Outcome-keyed placeholders:
%   H1 confirmed: "Lost in Canonicalization: Why Translation Does Not Close the
%                  Low-Resource Gap in Knowledge Graph Population"
%   B2:           "Translate, Then Extract: The Low-Resource Knowledge Graph
%                  Population Gap Is a Machine Translation Problem"
%   B3:           "Translation Is Not a Shortcut: Measuring Where Knowledge Graph
%                  Population Fails for Low-Resource Languages"
\title{[TITLE AFTER RESULTS — see placeholder block above]}
\author{Anonymous Author(s)}
\institute{Anonymous Institution}

\maketitle

\begin{abstract}
% Abstract skeleton (hypothesis package section 3) — sentence slots:
% 1. Populating knowledge bases such as Wikidata from text in low-resource languages
%    is limited by an unmeasured question: whether failures come from models not
%    understanding the text, or from models failing to ground entities into the
%    canonical knowledge base.
% 2. We test the hypothesis that translation to English repairs relation
%    identification but not entity grounding\footnote{Per the gap-recovery criterion
%    pre-registered in Section~\ref{sec:instrument}.}, using [N] extraction systems
%    across eight Indic languages in two arms — direct extraction and
%    translate-then-extract — with component-level scoring.
% 3. To make this measurable, we adapt the manually annotated portion of an existing
%    fact-to-text alignment resource into an extraction benchmark, and validate the
%    adaptation with a native-speaker audit plus automatic cross-checking,
%    quantifying its miss rate.
% 4. [RESULT SENTENCE — written after experiments; measured direction + effect size.]
% 5. We release the benchmark, audit protocol, and evaluation harness, designed so
%    new languages can be added without bespoke annotation.
\end{abstract}

\section{Introduction: The Hypothesis and Why It Matters}\label{sec:intro}
% First paragraph = the hypothesis, verbatim from the pre-registration prose form.
% No pipeline advertising. Banned vocabulary list: PREREGISTRATION.md section 9.

\section{Related Work}\label{sec:related}
% Assembled from the verified-source lists in the Maestro reports; every entry
% re-verified at writing time. \citeneed{XAlign WWW 2022} \citeneed{REDFM ACL 2023}
% \citeneed{Better To Ask in English} \citeneed{fertility literature}

\section{The Instrument: Benchmark and Its Validation}\label{sec:instrument}
% Reversed-XAlign benchmark; three-way triple classification; component metrics;
% the C2 native-speaker audit AS instrument validation (audit protocol = artifact);
% pre-registration summary + pointer to the frozen file (anonymized repo link).

\section{The Test: Direct vs.\ Translate-then-Extract}\label{sec:test}
% GRR results per language x system; decision-table verdicts; H1/B2/B3 branch outcome;
% English-ceiling pairing note + synthetic-parallel robustness probe result.

\section{Secondary Analyses}\label{sec:secondary}
% C3 (if day-8 gate passed: Hindi anchor + 7-language transfer test), C4/C5-fragment
% (exploratory, labeled), error taxonomy incl.\ hallucination structure (C6 dimension).

\section{Implications and Limitations}\label{sec:implications}
% For practitioners (translate first?) and benchmark builders (gold-reuse reliability).
% Limitations: single-family breadth, gold reuse caveats, audit on one language.

% \section*{AI-Assistance Disclosure} — Springer Nature policy statement, drafted day 13+.

\bibliographystyle{splncs04}
\bibliography{references} % references.bib starts EMPTY — entries only after verification

\end{document}
